\documentclass[UTF8,a4paper,12pt,openany,fontset=windows]{ctexrep}

\usepackage[left=3.0cm,right=2.5cm,top=2.6cm,bottom=2.6cm,headheight=15pt]{geometry}
\usepackage{fontspec}
\usepackage{xeCJK}
\usepackage{graphicx}
\usepackage{booktabs}
\usepackage{tabularx}
\usepackage{longtable}
\usepackage{array}
\usepackage{multirow}
\usepackage{float}
\usepackage{caption}
\usepackage{subcaption}
\usepackage{setspace}
\usepackage{fancyhdr}
\usepackage{enumitem}
\usepackage{amsmath}
\usepackage{amssymb}
\usepackage{hyperref}
\usepackage{xcolor}
\usepackage{microtype}

\setmainfont{Times New Roman}
\setsansfont{Arial}
\setmonofont{Consolas}
\setCJKmainfont{SimSun}
\setCJKsansfont{SimHei}
\setCJKmonofont{FangSong}

\definecolor{thesisblack}{RGB}{0,0,0}
\color{thesisblack}
\hypersetup{
  colorlinks=false,
  pdfborder={0 0 0},
  pdftitle={慧医数字医疗应用系统实习设计报告},
  pdfauthor={软件工程2302班02组}
}

\ctexset{
  chapter={
    format=\centering\heiti\zihao{-2}\bfseries\color{thesisblack},
    name={第,章},
    number=\arabic{chapter},
    beforeskip=8pt,
    afterskip=24pt,
    fixskip=true
  },
  section={
    format=\heiti\zihao{3}\bfseries\color{thesisblack},
    beforeskip=18pt,
    afterskip=10pt
  },
  subsection={
    format=\heiti\zihao{4}\bfseries\color{thesisblack},
    beforeskip=12pt,
    afterskip=6pt
  },
  contentsname={目\quad 录}
}

\setstretch{1.5}
\setlength{\emergencystretch}{3em}
\setlength{\parindent}{2em}
\setlength{\parskip}{0pt}
\setlist{nosep,leftmargin=2em}
\captionsetup{font=small,labelfont=bf,textfont=normalfont,labelsep=quad}
\renewcommand{\arraystretch}{1.35}
\IfFileExists{图片/01-用例图.png}{%
  \graphicspath{{图片/}}%
}{%
  \graphicspath{{../../最终交付文档/图片/}}%
}

\pagestyle{fancy}
\fancyhf{}
\fancyhead[C]{\songti\zihao{5}慧医数字医疗应用系统实习设计报告}
\fancyfoot[C]{\zihao{5}\thepage}
\renewcommand{\headrulewidth}{0.4pt}
\renewcommand{\footrulewidth}{0pt}
\fancypagestyle{plain}{
  \fancyhf{}
  \fancyhead[C]{\songti\zihao{5}慧医数字医疗应用系统实习设计报告}
  \fancyfoot[C]{\zihao{5}\thepage}
  \renewcommand{\headrulewidth}{0.4pt}
}

\newcolumntype{Y}{>{\centering\arraybackslash}X}
\newcolumntype{L}{>{\raggedright\arraybackslash}X}
\newcommand{\thesisfigure}[3][0.96\textwidth]{%
  \begin{figure}[H]
    \centering
    \includegraphics[width=#1,height=0.72\textheight,keepaspectratio]{#2}
    \caption{#3}
  \end{figure}
}
\newcommand{\placeholder}[1]{\colorbox{yellow}{\strut\songti\zihao{5}【待填写：#1】}}
\newcommand{\codefont}{\ttfamily\footnotesize}

\begin{document}
\songti\zihao{-4}
\hypersetup{pageanchor=false}

% 封面
\begin{titlepage}
  \thispagestyle{empty}
  \begin{center}
    \vspace*{1.3cm}
    {\heiti\zihao{1}\bfseries 实习设计报告\par}
    \vspace{1.8cm}
    {\heiti\zihao{2}\bfseries 慧医数字医疗应用系统\par}
    \vspace{0.8cm}
    {\songti\zihao{3}软件工程综合实训项目\par}
    \vfill
    \begin{minipage}{0.78\textwidth}
      \zihao{4}
      \noindent\heiti 设计题目：\songti 慧医数字医疗应用系统设计与实现\par\vspace{0.55em}
      \noindent\heiti 班级：\songti 软件工程2302\par\vspace{0.55em}
      \noindent\heiti 小组名称：\songti 02组\par\vspace{0.55em}
      \noindent\heiti 小组组长：\songti 李梦杰\par\vspace{0.55em}
      \noindent\heiti 小组成员：\songti 曾子唯、詹宁康、周林桂、侬常旭\par\vspace{0.55em}
      \noindent\heiti 指导教师：\songti 万华军\par\vspace{0.55em}
      \noindent\heiti 完成日期：\songti 2026年7月17日\par
    \end{minipage}
    \vfill
    {\songti\zihao{4}2026年7月\par}
  \end{center}
\end{titlepage}

\hypersetup{pageanchor=true}
\pagenumbering{Roman}
\chapter*{摘\quad 要}
\addcontentsline{toc}{chapter}{摘要}

慧医数字医疗应用系统是一套面向医疗基础数据维护与业务查询的前后端分离应用。系统围绕管理员与医生两类角色，完成登录认证、基于角色的访问控制、医生信息、医药公司、销售地点、城市、药品、医保政策、公司政策和必备材料等业务模块，并提供首页统计、药品图片安全上传、销售地点地图展示、健康检查和持续集成部署能力。

本报告依据仓库最新代码、数据库脚本、测试报告、Docker 与 Nginx 配置、CodeArts 流水线资料、GitHub 协作记录以及最终交付图片编写。报告采用需求驱动与分层设计方法，依次完成需求分析、概要设计、详细设计、数据库设计、关键实现、测试验证、构建部署和项目管理说明。系统前端采用 Vue 3、Vue Router、Vuex、Element Plus 与 Vite，后端采用 Spring Boot 2.7.18、Spring Security、MyBatis、MySQL 8 和 Redis 7，生产环境通过 Docker Compose 编排，并由 Nginx 提供同源反向代理。

截至报告整理时，仓库代码级自动化测试共 215 项并全部通过，其中后端 84 项、前端 123 项、Python 8 项。后端行覆盖率和方法覆盖率均为 100\%，前端语句覆盖率和行覆盖率均为 100\%。历史容器环境还完成 57 项 API 回归，证明构建、数据初始化、接口和部署链路能够形成闭环。

\noindent\textbf{关键词：}数字医疗；前后端分离；Spring Boot；Vue；持续集成；软件测试

\clearpage
\tableofcontents
\clearpage
\listoffigures
\addcontentsline{toc}{chapter}{插图目录}
\clearpage
\listoftables
\addcontentsline{toc}{chapter}{表格目录}

\clearpage
\pagenumbering{arabic}

\chapter{绪论}

\section{项目背景}

医疗机构及医药相关组织在日常管理中需要维护医生、药品、政策、材料、企业和销售地点等多类基础数据。传统分散式表格管理容易产生数据重复、权限边界模糊、查询效率低和变更难以追踪等问题。为解决上述问题，本项目建设统一的数字医疗应用系统，通过浏览器提供一致的操作入口，并通过服务端集中实现认证、授权、数据校验、审计和持久化。

\section{设计目标}

项目目标可概括为四个方面：第一，形成覆盖核心业务对象的完整信息管理能力；第二，以管理员和医生角色建立清晰的最小权限边界；第三，保证系统能够在本地、容器和云端环境中重复构建与部署；第四，通过自动化测试、代码检查和持续集成提升交付质量并保留可复核证据。

\section{报告范围与研究方法}

本报告覆盖项目需求、架构、模块、数据库、接口、安全、测试、部署和协作全过程。研究方法包括需求追踪、代码静态审阅、配置审阅、数据库脚本分析、UML 建模、测试结果复核和平台证据归档。对于不能从仓库事实确认的生产域名等信息，报告保留明确的待填写标记，不将假设写成已上线事实。

\section{项目成员与分工}

\begin{table}[H]
  \centering
  \caption{项目小组信息}
  \begin{tabularx}{\textwidth}{p{3.2cm}L}
    \toprule
    项目项 & 内容 \\
    \midrule
    班级 & 软件工程2302 \\
    小组名称 & 02组 \\
    小组组长 & 李梦杰 \\
    小组成员 & 曾子唯、詹宁康、周林桂、侬常旭 \\
    指导教师 & 万华军 \\
    完成时间 & 2026年7月17日 \\
    \bottomrule
  \end{tabularx}
\end{table}

\chapter{需求分析}

\section{用户角色分析}

系统包含管理员与医生两类主要角色。管理员具有业务基础数据的维护权限，负责新增、修改、删除和查询；医生以业务查询和使用为主，不应取得系统级维护权限。所有受保护请求必须经过会话校验和角色授权，未登录用户只能访问登录、会话探测和必要的健康检查入口。

\begin{table}[H]
  \centering
  \caption{角色与权限需求}
  \begin{tabularx}{\textwidth}{p{2.8cm}p{3.5cm}L}
    \toprule
    角色 & 权限范围 & 主要职责 \\
    \midrule
    管理员（ROLE\_1） & 全量业务维护 & 管理账号、医生、药品、企业、政策、城市、销售地点、材料和图片 \\
    医生（ROLE\_2） & 授权范围内查询 & 查看首页统计和业务信息，使用地图与资讯，不执行越权维护 \\
    未登录用户 & 公开入口 & 登录、会话状态探测和网关允许的健康检查 \\
    \bottomrule
  \end{tabularx}
\end{table}

\section{功能需求}

\begin{longtable}{p{1.4cm}p{3.1cm}p{4.6cm}p{4.6cm}}
  \caption{主要功能需求}\label{tab:functional-requirements}\\
  \toprule
  编号 & 模块 & 功能说明 & 验收要点 \\
  \midrule
  \endfirsthead
  \toprule
  编号 & 模块 & 功能说明 & 验收要点 \\
  \midrule
  \endhead
  F01 & 登录认证 & 校验账号、密码与账号状态，建立服务端会话 & 正确凭据登录；错误凭据拒绝；Cookie 为 httpOnly \\
  F02 & 权限控制 & 根据真实角色返回权限树并控制接口访问 & 管理员与医生权限隔离；越权请求返回拒绝 \\
  F03 & 首页统计 & 聚合医生、药品、城市等数量与资讯 & 数据可视化正确；业务变更后缓存失效 \\
  F04 & 医生管理 & 医生信息分页、查询、新增、修改、删除与密码重置 & 唯一性、状态和权限校验正确 \\
  F05 & 医药公司 & 企业基础信息与公司政策维护 & 分页和增删改查完整 \\
  F06 & 销售地点 & 销售地点与经纬度维护，地图展示与选点 & 坐标可保存、回显并在地图中定位 \\
  F07 & 城市管理 & 城市信息分页和维护 & 数据校验与关联使用正常 \\
  F08 & 药品管理 & 药品信息、价格、公司关联与图片维护 & 图片路径安全，列表与详情一致 \\
  F09 & 政策管理 & 医保政策与公司政策维护 & 富文本或长文本完整保存和查询 \\
  F10 & 材料管理 & 必备材料信息维护 & 数据增删改查和权限控制正确 \\
  F11 & 健康检查 & 汇总应用、MySQL 和 Redis 状态 & 容器与流水线可以据此判定部署成功 \\
  \bottomrule
\end{longtable}

\section{非功能需求}

\begin{table}[H]
  \centering
  \caption{非功能需求}
  \begin{tabularx}{\textwidth}{p{2.6cm}L p{4.0cm}}
    \toprule
    类别 & 需求描述 & 评价方式 \\
    \midrule
    安全性 & 密码摘要存储；会话不可被脚本读取；文件名和路径受控；生产后端端口不直接暴露公网 & 安全测试、配置检查、越权用例 \\
    可靠性 & 应用异常统一处理；数据库和 Redis 状态可观测；部署失败可回滚 & 自动化测试、健康检查、流水线记录 \\
    可维护性 & 前后端分层、模块边界清楚、配置外置、数据库迁移可追踪 & 代码结构和文档审查 \\
    性能 & 列表分页；首页聚合结果缓存；静态资源由 Nginx 提供 & 接口耗时与历史 API 回归数据 \\
    可部署性 & 支持 JDK 17、Node 18、Docker Compose 与 CodeArts 流水线 & 构建、容器启动和云端部署验证 \\
    可测试性 & 核心模块具有单元或组件测试，覆盖率可量化 & Maven、Vitest/V8、Python 测试报告 \\
    \bottomrule
  \end{tabularx}
\end{table}

\section{业务边界与约束}

系统以 MySQL 数据库 \texttt{medicine} 为权威业务数据源，以 Redis 保存会话与缓存。前端不得直接连接数据库或 Redis；所有写操作通过后端接口完成。生产环境通过 Nginx 暴露 80/443 端口，后端 8082 仅在主机内部或容器网络中可达。药品图片只能写入指定上传目录，并限制扩展名、MIME 类型、大小和规范化路径。

\thesisfigure{01-用例图.png}{系统用例图}

\chapter{概要设计}

\section{设计思想}

系统采用“分层解耦、权限前置、配置外置、测试护航、容器交付”的总体思想。浏览器端负责交互与展示，Nginx 负责统一入口和反向代理，Spring Boot 负责业务规则、安全和数据访问，MySQL 负责持久化，Redis 负责会话和热点缓存。所有业务模块遵循一致的控制器、服务、数据访问和领域对象组织方式，从而降低模块之间的直接依赖。

\section{实现方法}

前端使用 Vue 3 单文件组件组织页面，以 Vue Router 实现动态路由，以 Vuex 保存登录态和权限，以 Axios 调用同源 \texttt{/api} 接口。后端使用 Spring MVC 暴露 REST 接口，Spring Security 建立认证授权链，MyBatis 完成 SQL 映射，统一异常处理器返回一致响应结构。数据库变更通过基线脚本和 V2、V3、V4 迁移文件演进。构建产物由 Docker Compose 编排，并通过 CodeArts 完成代码检查、构建、制品、部署和接口健康验证。

\section{系统总体架构}

\begin{table}[H]
  \centering
  \caption{系统分层架构}
  \begin{tabularx}{\textwidth}{p{2.6cm}p{4.0cm}L}
    \toprule
    层次 & 主要技术 & 职责 \\
    \midrule
    表示层 & Vue 3、Element Plus、ECharts、地图组件 & 页面交互、表单校验、数据展示和路由导航 \\
    接入层 & Nginx & 静态资源、同源代理、TLS 终止、请求头转发 \\
    应用层 & Spring Boot、Spring MVC、Spring Security & 接口编排、认证授权、统一响应与异常处理 \\
    领域与服务层 & Service、事件、缓存维护任务 & 业务规则、事务、数据变更事件与缓存失效 \\
    数据访问层 & MyBatis & SQL 映射、分页查询和实体持久化 \\
    基础设施层 & MySQL 8、Redis 7、Docker Compose & 持久数据、会话缓存、容器运行和健康检查 \\
    \bottomrule
  \end{tabularx}
\end{table}

\section{主要模块及模块关系}

认证模块位于所有受保护模块之前，向前端返回会话状态和权限树；业务模块共享统一的分页、响应和异常处理设施；医生、药品、企业、城市、政策、材料和销售地点等模块通过数据库外键或业务标识建立关系；首页面板读取多个业务模块的聚合数据，并通过业务变更事件触发缓存失效；图片模块为药品模块提供受控文件存储；部署模块为全部运行组件提供环境变量、Secret、健康检查和网络边界。

\thesisfigure{02-领域类图.png}{领域类图}
\thesisfigure{06-组件图.png}{系统组件图}
\thesisfigure{07-包图.png}{系统包图}
\thesisfigure{08-部署图.png}{系统部署图}

\section{技术选型}

\begin{table}[H]
  \centering
  \caption{主要技术选型}
  \begin{tabularx}{\textwidth}{p{2.8cm}p{4.2cm}L}
    \toprule
    类别 & 技术与版本 & 选型说明 \\
    \midrule
    前端框架 & Vue 3.5、Vue Router 4、Vuex 4 & 组件化开发、动态路由和集中状态管理 \\
    构建与测试 & Vite 6、Vitest、V8 Coverage & 快速构建、组件测试和覆盖率统计 \\
    后端框架 & Spring Boot 2.7.18、JDK 17 & 成熟的 Web、配置、测试和健康检查生态 \\
    安全与持久化 & Spring Security、MyBatis & 统一安全链与可控 SQL 映射 \\
    数据与缓存 & MySQL 8、Redis 7 & 关系数据持久化、会话和热点数据缓存 \\
    交付平台 & Docker Compose、Nginx、CodeArts & 环境一致性、反向代理和持续交付 \\
    \bottomrule
  \end{tabularx}
\end{table}

\chapter{详细设计}

\section{后端类与接口设计}

后端按业务域组织控制器、服务、数据访问对象和实体。控制器只处理协议层参数与响应，服务层承担业务校验和事务，Mapper 负责数据库访问。公共模块提供统一响应、异常处理、分页支持、请求耗时拦截器和配置对象，避免各业务模块重复实现基础逻辑。

\thesisfigure{03-后端设计类图.png}{后端设计类图}
\thesisfigure{05-对象图.png}{典型运行时对象图}

\section{登录认证与规范化 RBAC 设计}

登录时后端校验账号状态和 BCrypt 密码摘要，成功后生成随机不透明令牌，仅将令牌的 SHA-256 摘要键与最小会话身份写入 Redis，并通过 httpOnly Cookie 交付浏览器。后续请求由安全过滤链提取 Cookie、校验 Redis 会话，再以会话中的可信账号编号为入口重新加载当前授权。原始令牌不写入 Redis，也不在登录响应体中暴露。

V5 迁移将原来的“账号字符串角色 + 菜单映射”升级为规范化的账号--角色--权限模型。核心表映射见表\ref{tab:rbac-tables}；一个账号可以关联多个角色，菜单权限和接口动作权限使用同一权限目录，但通过类型字段严格区分。旧字段与旧关系表只作为数据回填和旧版本回滚依据，不再作为新授权查询入口。

\begin{table}[H]
  \centering
  \caption{规范化 RBAC 核心表映射}
  \label{tab:rbac-tables}
  \begin{tabularx}{\textwidth}{p{3.2cm}p{5.1cm}L}
    \toprule
    表 & 关键字段 & 职责 \\
    \midrule
    \texttt{account} & \texttt{id}、\texttt{status}、\texttt{utype} & 认证主体与启停状态；\texttt{utype} 仅为兼容字段 \\
    \texttt{rbac\_role} & \texttt{id}、\texttt{code}、\texttt{enabled} & 角色主数据，当前启用 \texttt{ADMIN} 与 \texttt{DOCTOR} \\
    \texttt{account\_role} & \texttt{account\_id}、\texttt{role\_id}、\texttt{is\_primary} & 账号与角色的多对多关系，并约束每个账号至多一个主角色 \\
    \texttt{permission} & \texttt{code}、\texttt{permission\_type}、\texttt{enabled}、\texttt{sort\_order} & 统一权限目录；\texttt{MENU} 表示动态菜单，\texttt{ACTION} 表示接口动作 \\
    \texttt{rbac\_role\_permission} & \texttt{role\_id}、\texttt{permission\_id} & 角色与菜单/API 权限的多对多授权关系 \\
    \texttt{role\_permission} & \texttt{roleName}、\texttt{per\_id} & V5 回填来源及旧版本回滚数据，不参与新授权链路 \\
    \bottomrule
  \end{tabularx}
\end{table}

运行时授权链路为“\texttt{accountId} $\rightarrow$ \texttt{account\_role} $\rightarrow$ 启用的 \texttt{rbac\_role} $\rightarrow$ \texttt{rbac\_role\_permission} $\rightarrow$ 启用的 \texttt{permission}”。\texttt{TokenAuthenticationFilter} 从会话主体取得 \texttt{accountId}，逐次请求联查账号状态、角色和权限，将形如 \texttt{ROLE\_ADMIN}、\texttt{ROLE\_DOCTOR} 的角色权限及 \texttt{resource:verb} 动作权限写入 Spring Security 的 \texttt{GrantedAuthority}。账号不存在、账号被禁用或没有启用角色时不建立认证上下文；缺少具体动作权限时按默认拒绝处理。因此角色撤销、权限调整和账号禁用在下一次请求即可生效，不受 Redis 会话剩余时间影响。

前后端按“展示控制”和“安全控制”分层：\texttt{/api/permissions} 只根据认证主体的 \texttt{accountId} 查询多角色权限并集，返回 \texttt{MENU} 菜单树、\texttt{permissionCodes} 和 \texttt{roles}，客户端即使传入 \texttt{roleName} 也不会参与授权。Vuex、动态路由和 \texttt{\$can()} 只负责菜单、按钮与交互提示；后端控制器通过 \texttt{@PreAuthorize(hasAuthority(...))} 校验 \texttt{ACTION} 权限，才是不可绕过的最终安全边界。管理员获得全量动作权限，医生保持业务只读，不因前端元素是否可见而改变服务端结论。

\thesisfigure{09-登录活动图.png}{登录活动图}
\thesisfigure{12-账号状态图.png}{账号状态图}
\thesisfigure{14-登录授权顺序图.png}{登录授权顺序图}
\thesisfigure{17-重置密码顺序图.png}{重置密码顺序图}

\section{业务数据管理设计}

业务管理页面采用统一的“条件查询、分页列表、新增、编辑、删除”交互模式。前端分页组件统一页码和每页数量，API 模块封装路径与参数，后端控制器将请求交由对应服务处理。新增和修改操作先做必填、格式和唯一性校验，再执行事务写入；删除操作检查权限与关联约束，成功后发布业务数据变更事件。

\thesisfigure{10-医生管理活动图.png}{医生管理活动图}
\thesisfigure{15-新增医生顺序图.png}{新增医生顺序图}

\section{药品图片上传设计}

图片上传由后端统一控制：只允许配置的图片类型和扩展名，限制文件大小，对文件名进行随机化处理，并将目标路径规范化后验证其仍位于上传根目录内。数据库只保存可访问的相对路径。Nginx 的 \texttt{/image/} 路径映射到受控目录，并设置类型限制、\texttt{nosniff} 和缓存策略。

\thesisfigure{11-药品图片上传活动图.png}{药品图片上传活动图}
\thesisfigure{16-药品保存上传顺序图.png}{药品保存与图片上传顺序图}

\section{首页面板与缓存设计}

首页面板需要聚合多个业务表。系统将统计数量、分布和资讯结果缓存在 Redis 中，降低重复查询成本。业务服务完成数据变更后发布事件，监听器统一删除相关缓存；定时维护任务作为补充机制，防止长时间保留陈旧数据。缓存不可用时，服务仍可从数据库读取，保证核心业务可降级运行。

\thesisfigure{13-资讯材料状态图.png}{资讯与材料状态图}
\thesisfigure{18-首页面板顺序图.png}{首页面板顺序图}

\chapter{数据库设计}

\section{数据模型}

数据库使用 \texttt{medicine} 模式，以账号、医生、企业、城市、销售地点、药品、医保政策、公司政策和材料等表保存核心数据。基线脚本定义初始结构，迁移脚本逐步补充业务字段、索引、经纬度坐标、审计记录和规范化 RBAC 关系。执行 V2--V5 后共包含 23 张业务及关联表，其中 V5 在原 20 张表基础上新增 \texttt{rbac\_role}、\texttt{account\_role} 和 \texttt{rbac\_role\_permission} 3 张关系表。

\thesisfigure{04-实体关系图ER.png}{数据库实体关系图}

\section{核心数据表}

\begin{table}[H]
  \centering
  \caption{核心数据表说明}
  \begin{tabularx}{\textwidth}{p{3.3cm}p{3.7cm}L}
    \toprule
    表或领域 & 关键字段 & 设计说明 \\
    \midrule
    账号 & id、username、password、status & 保存认证主体、BCrypt 摘要密码和启停状态 \\
    角色与账号关系 & rbac\_role、account\_role & 保存稳定角色编码、启停状态和账号多角色关系 \\
    权限与角色关系 & permission、rbac\_role\_permission & 统一保存 MENU/ACTION 权限码及角色授权 \\
    医生 & id、accountId、name、hospital & 连接登录账号与医生业务档案 \\
    医药公司 & id、name、address、contact & 作为药品和公司政策的业务主体 \\
    城市与销售地点 & cityId、saleId、longitude、latitude & 支撑行政区域、销售网点和地图定位 \\
    药品 & id、name、companyId、price、image & 保存药品属性、公司关联和图片相对路径 \\
    医保与公司政策 & id、title、content、publishTime & 保存政策文本及发布信息 \\
    必备材料 & id、name、description、status & 保存业务所需材料和有效状态 \\
    \bottomrule
  \end{tabularx}
\end{table}

\section{数据库演进与初始化}

\begin{table}[H]
  \centering
  \caption{数据库脚本演进}
  \begin{tabularx}{\textwidth}{p{3.2cm}p{3.0cm}L}
    \toprule
    脚本 & 阶段 & 主要作用 \\
    \midrule
    \texttt{medical.sql} & 基线 & 创建初始业务表、约束和基础数据 \\
    V2 迁移 & 需求对齐 & 增加账号状态、销售地点坐标、资讯、数据质量和密码重置审计等结构 \\
    V3 迁移 & 审计生命周期 & 将密码重置审计外键调整为随账号删除级联清理 \\
    V4 迁移 & 地图菜单 & 增加销售地点地图菜单并补充旧角色菜单映射 \\
    V5 迁移 & 规范化 RBAC & 新增账号--角色--权限关系、MENU/ACTION 权限码及兼容数据回填 \\
    \bottomrule
  \end{tabularx}
\end{table}

初始化顺序固定为 \texttt{medical.sql} $\rightarrow$ \texttt{V2\_\_requirements\_alignment.sql} $\rightarrow$ \texttt{V3\_\_password\_audit\_delete\_policy.sql} $\rightarrow$ \texttt{V4\_\_sale\_map\_menu.sql} $\rightarrow$ \texttt{V5\_\_normalized\_rbac.sql}。脚本执行完毕后应校验 23 张表、关键索引与外键、权限码唯一性、账号主角色唯一性，以及管理员全量权限和医生无写权限，最后再启动应用。生产环境执行迁移前必须备份数据库，并保证脚本具备可追踪版本，不允许直接修改已经发布的历史迁移文件。

\chapter{系统实现}

\section{核心代码实现}

本节仅选取能够说明关键设计的短代码片段，不复制完整类或长函数。所有代码均在三线表中展示，完整实现以仓库源代码为准。

\begin{table}[H]
  \centering
  \caption{核心代码：公共入口与认证边界}
  \begin{tabularx}{\textwidth}{p{1.2cm}L}
    \toprule
    行 & 代码 \\
    \midrule
    1 & \codefont .antMatchers(HttpMethod.OPTIONS, "/**").permitAll() \\
    2 & \codefont .antMatchers(HttpMethod.GET, "/image/**").permitAll() \\
    3 & \codefont .antMatchers("/api/login", "/api/session", \\
    4 & \codefont \quad "/actuator/health", "/actuator/info").permitAll() \\
    \bottomrule
  \end{tabularx}
\end{table}

\begin{table}[H]
  \centering
  \caption{核心代码：安全会话 Cookie}
  \begin{tabularx}{\textwidth}{p{1.2cm}L}
    \toprule
    行 & 代码 \\
    \midrule
    1 & \codefont ResponseCookie.ResponseCookieBuilder builder = \\
    2 & \codefont \quad ResponseCookie.from(cookieProperties.getName(), token) \\
    3 & \codefont \quad .httpOnly(true).secure(cookieProperties.isSecure()) \\
    4 & \codefont \quad .sameSite(cookieProperties.getSameSite()) \\
    5 & \codefont \quad .path(cookieProperties.getPath()); \\
    \bottomrule
  \end{tabularx}
\end{table}

\begin{table}[H]
  \centering
  \caption{核心代码：业务变更触发缓存失效}
  \begin{tabularx}{\textwidth}{p{1.2cm}L}
    \toprule
    行 & 代码 \\
    \midrule
    1 & \codefont @AfterReturning("businessMutation()") \\
    2 & \codefont public void publishDataChangedEvent(JoinPoint joinPoint) \{ \\
    3 & \codefont \quad eventPublisher.publishEvent( \\
    4 & \codefont \quad\quad new BusinessDataChangedEvent( \\
    5 & \codefont \quad\quad joinPoint.getSignature().toShortString())); \\
    6 & \codefont \} \\
    \bottomrule
  \end{tabularx}
\end{table}

\begin{table}[H]
  \centering
  \caption{核心代码：前端统一 API 请求}
  \begin{tabularx}{\textwidth}{p{1.2cm}L}
    \toprule
    行 & 代码 \\
    \midrule
    1 & \codefont export function getDoctorInfo(pn, size, keyword = '') \{ \\
    2 & \codefont \quad return request(\{ url: '/doctors', method: 'GET', \\
    3 & \codefont \quad\quad params: \{ pn, size, keyword \} \}) \\
    4 & \codefont \quad\quad .then((res) $=>$ judgeQueryResult(res)); \\
    5 & \codefont \} \\
    \bottomrule
  \end{tabularx}
\end{table}

\begin{table}[H]
  \centering
  \caption{核心代码：Nginx 同源代理}
  \begin{tabularx}{\textwidth}{p{1.2cm}L}
    \toprule
    行 & 代码 \\
    \midrule
    1 & \codefont location /api/ \{ \\
    2 & \codefont \quad proxy\_pass \$medicine\_backend\_upstream\$request\_uri; \\
    3 & \codefont \quad proxy\_http\_version 1.1; \\
    4 & \codefont \quad proxy\_set\_header Origin ""; \\
    5 & \codefont \} \\
    \bottomrule
  \end{tabularx}
\end{table}

\begin{table}[H]
  \centering
  \caption{核心代码：Docker Compose 健康检查}
  \begin{tabularx}{\textwidth}{p{1.2cm}L}
    \toprule
    行 & 代码 \\
    \midrule
    1 & \codefont healthcheck: \\
    2 & \codefont \quad test: ["CMD-SHELL", "wget -q -O - \\
    3 & \codefont \quad\quad http://127.0.0.1:8082/actuator/health \\
    4 & \codefont \quad\quad | grep -q '"status":"UP"'"] \\
    5 & \codefont \quad interval: 10s \\
    \bottomrule
  \end{tabularx}
\end{table}

\section{系统运行效果}

\thesisfigure{项目效果登录页面.png}{系统登录页面}
\thesisfigure{项目效果首页.png}{系统首页统计面板}
\thesisfigure{项目效果医生信息管理.png}{医生信息管理页面}
\thesisfigure{项目效果医药信息管理.png}{药品信息管理页面}
\thesisfigure{项目效果药品信息管理.png}{药品维护操作页面}
\thesisfigure{项目效果医药公司政策管理.png}{医药公司政策管理页面}
\thesisfigure{项目效果医保政策管理.png}{医保政策管理页面}
\thesisfigure{项目效果城市信息管理.png}{城市信息管理页面}
\thesisfigure{项目效果销售地点管理.png}{销售地点管理页面}
\thesisfigure{项目效果销售地点地图管理.png}{销售地点地图管理页面}
\thesisfigure{项目效果新增地点1.png}{新增销售地点表单}
\thesisfigure{项目效果新增地点2.png}{新增地点地图选点}
\thesisfigure{项目效果必备材料管理.png}{必备材料管理页面}
\thesisfigure{项目效果增删改查.png}{通用增删改查操作效果}

\chapter{系统测试}

\section{测试策略与环境}

测试采用分层策略：后端使用 JUnit、Mockito、Spring Boot Test 与 JaCoCo；前端使用 Vitest、Vue Test Utils 与 V8 Coverage；Python 测试验证 API 回归脚本本身；历史部署环境使用 57 项黑盒 API 用例验证容器、数据库、Redis 与 Nginx 链路。测试环境以 JDK 17、Maven 3.8 以上、Node 18 以上、MySQL 8、Redis 7 和 Docker Compose v2 为基线。

\begin{table}[H]
  \centering
  \caption{代码级自动化测试结果}
  \begin{tabularx}{\textwidth}{p{3.1cm}p{2.5cm}p{2.5cm}L}
    \toprule
    测试层次 & 通过数 & 总数 & 结论 \\
    \midrule
    后端 JUnit 测试 & 100 & 100 & 全部通过 \\
    前端 Vitest 测试 & 133 & 133 & 全部通过 \\
    Python 测试 & 8 & 8 & 全部通过 \\
    合计 & 241 & 241 & 通过率 100\% \\
    \bottomrule
  \end{tabularx}
\end{table}

\section{覆盖率结果}

\begin{table}[H]
  \centering
  \caption{自动化测试覆盖率}
  \begin{tabularx}{\textwidth}{p{3.0cm}p{3.0cm}p{2.6cm}L}
    \toprule
    对象 & 指标 & 覆盖率 & 说明 \\
    \midrule
    后端 & 行覆盖率 & 100\%（935/935） & 达到行覆盖率目标 \\
    后端 & 方法覆盖率 & 100\% & 全部方法被执行 \\
    后端 & 指令覆盖率 & 99.74\% & 少量编译级指令未覆盖 \\
    后端 & 分支覆盖率 & 87.39\% & 复杂异常组合仍可继续补充 \\
    前端 & 语句覆盖率 & 100\% & 全部语句被执行 \\
    前端 & 行覆盖率 & 100\% & 全部代码行被执行 \\
    前端 & 分支覆盖率 & 95.64\% & 少量防御分支未命中 \\
    前端 & 函数覆盖率 & 63.88\% & 视图回调和框架绑定仍有提升空间 \\
    \bottomrule
  \end{tabularx}
\end{table}

\section{历史部署环境 API 回归}

2026年7月13日的容器化环境完成 57 项 API 回归并全部通过，平均响应时间约 309.9 ms，P95 为 550 ms，最大值为 639 ms。该结果用于证明当时版本的部署链路完整性，与 2026年7月15日整理的 215 项代码级测试属于不同测试批次，不将历史运行结果冒充本次重新执行。

\begin{table}[H]
  \centering
  \caption{历史 API 回归结果}
  \begin{tabularx}{\textwidth}{p{3.1cm}p{2.6cm}p{2.6cm}L}
    \toprule
    指标 & 数值 & 状态 & 说明 \\
    \midrule
    API 用例 & 57/57 & 通过 & 覆盖认证、权限、业务查询和写操作 \\
    平均响应时间 & 309.9 ms & 正常 & 历史容器测试环境统计 \\
    P95 响应时间 & 550 ms & 正常 & 95\% 请求不超过该值 \\
    最大响应时间 & 639 ms & 正常 & 未发现超时或失败请求 \\
    MySQL/Redis & UP/UP & 通过 & 健康检查确认依赖可用 \\
    \bottomrule
  \end{tabularx}
\end{table}

\section{代码检查问题整改}

2026年7月15日依据华为云 CodeArts 代码扫描报告与代码评审清单，对致命与严重问题集中整改，每项修复均附回归测试，并以小粒度提交同步至 GitHub 与 CodeArts 双仓库。整改后后端 JUnit 测试由 84 项增至 100 项，前端 Vitest 测试由 123 项增至 133 项。

\begin{table}[H]
  \centering
  \caption{代码检查问题整改清单}
  \begin{tabularx}{\textwidth}{p{1.1cm}p{3.6cm}L}
    \toprule
    级别 & 问题 & 整改措施 \\
    \midrule
    致命 & 临时密码使用 \texttt{Math.random()}，非密码学安全随机数（G.OTH.01） & 改用 \texttt{SecureRandom} 生成，补充密码格式回归测试 \\
    严重 & 嵌套三元表达式（G.EXP.03） & 拆分为单层三元与中间变量 \\
    严重 & import 未按 com/org/java/javax 分组（G.FMT.03） & 6 个生产文件补空行分组与字母排序 \\
    缺陷 & 核心业务表硬删除致误删不可恢复 & doctor/drug/policy/material/sale 改 \texttt{deleted\_at} 软删除，查询过滤，关联保留可恢复 \\
    缺陷 & 数据改动无操作者追溯 & 新增 \texttt{create\_by}/\texttt{update\_by}/\texttt{deleted\_by} 审计列，由登录账号 id 写入 \\
    缺陷 & 业务异常统一返回 200，状态码语义缺失 & \texttt{GlobalExceptionHandler} 映射 400/401/403/404/409/500，EntryPoint 与 DeniedHandler 返回 401/403 \\
    缺陷 & \texttt{DrugService.update} 缺 saleIds 误清空销售关联 & 仅当请求显式携带 saleIds 才重建关联 \\
    缺陷 & 搜索每次按键发请求、翻页与搜索页码耦合 & 新增 debounce 工具，6 个管理页防抖 300\,ms，新关键字回到第一页 \\
    缺陷 & DoctorManage 修改提交携带 \texttt{pwd:undefined} & 移除密码字段，密码走独立重置流程 \\
    加固 & 无内容安全策略与定时备份 & nginx 下发 CSP；新增 mysqldump 一致性备份脚本与 systemd timer \\
    \bottomrule
  \end{tabularx}
\end{table}

\section{测试结论与改进}

当前测试能够证明核心代码、组件、接口脚本和历史部署链路均具备较高可信度。需要说明的是，行覆盖率 100\% 不等于所有业务组合均无缺陷；后续仍应增加浏览器端到端测试、多浏览器兼容、可访问性、压力、故障注入、安全扫描和生产环境回归，并将前端函数覆盖率与前后端分支覆盖率作为持续改进指标。

\thesisfigure{华为云测试用例.png}{CodeArts 测试用例管理}
\thesisfigure{华为云代码检查.png}{CodeArts 代码检查结果}

\chapter{构建、部署与持续集成}

\section{本地构建与容器运行}

后端在 \texttt{medical-backend} 目录执行 \texttt{mvn verify}，生成包含测试门禁的可执行 JAR；前端在 \texttt{medical-managerment-system} 目录执行 \texttt{npm run test:coverage} 和 \texttt{npm run build}，生成 \texttt{dist} 静态资源。根目录执行 \texttt{docker compose up -d --build} 可构建并启动 \texttt{medicine-backend} 与 \texttt{medicine-web} 两个服务。浏览器默认访问 \texttt{127.0.0.1:9092}。

\begin{table}[H]
  \centering
  \caption{构建与运行入口}
  \begin{tabularx}{\textwidth}{p{2.8cm}p{5.2cm}L}
    \toprule
    对象 & 命令或入口 & 结果 \\
    \midrule
    后端 & \texttt{mvn verify} & 执行测试、覆盖率门禁并生成 JAR \\
    前端测试 & \texttt{npm run test:coverage} & 执行 123 项测试并生成 V8 覆盖率 \\
    前端构建 & \texttt{npm run build} & 生成生产静态资源 \texttt{dist/} \\
    Compose & \texttt{docker compose up -d --build} & 构建并启动前后端容器 \\
    健康检查 & \texttt{/actuator/health} & 返回应用、MySQL 和 Redis 状态 \\
    \bottomrule
  \end{tabularx}
\end{table}

\section{端口转发与反向代理}

开发或演示环境对外暴露 Web/Nginx 的 9092 端口，Nginx 将 \texttt{/api/}、\texttt{/image/} 和 \texttt{/actuator/health} 分别转发至后端 API、图片目录和健康检查。Compose 将后端 8082 绑定至宿主机回环地址，便于本机诊断；生产环境不应将 8082 直接暴露公网，而应只开放 80/443，由外层 Nginx 或云网关统一接入。

\begin{table}[H]
  \centering
  \caption{端口与路径转发}
  \begin{tabularx}{\textwidth}{p{3.1cm}p{3.6cm}L}
    \toprule
    入口 & 目标 & 用途 \\
    \midrule
    9092 或 80/443 & \texttt{medicine-web:80} & 页面统一入口 \\
    \texttt{/api/} & \texttt{backend:8082/api/} & 业务 API \\
    \texttt{/image/} & 后端受控上传目录 & 药品图片访问 \\
    \texttt{/actuator/health} & \texttt{backend:8082} & 健康检查 \\
    \texttt{127.0.0.1:8082} & Spring Boot & 仅限本机诊断 \\
    \bottomrule
  \end{tabularx}
\end{table}

\section{子域名与 HTTPS 配置}

本项目正式子域名已确认为 \texttt{sx.weavecan.com}，通过 DNS 指向 ECS 并经 SSH 隧道转发至本地 Docker 部署，TLS 证书已配置。Nginx \texttt{server\_name}、后端 CORS 白名单与 Cookie 安全策略均已按该域名同步。正式访问地址为：\url{https://sx.weavecan.com/}。

\begin{table}[H]
  \centering
  \caption{生产域名上线核对项}
  \begin{tabularx}{\textwidth}{p{3.5cm}p{4.2cm}L}
    \toprule
    配置项 & 建议值 & 说明 \\
    \midrule
    DNS & \texttt{sx.weavecan.com} 指向 ECS/负载均衡 & 根据公网网络填写 A、AAAA 或 CNAME \\
    Nginx \texttt{server\_name} & \texttt{sx.weavecan.com} & 替换演示配置中的主机名 \\
    CORS & \texttt{https://sx.weavecan.com} & Cookie 模式不能与通配符同时使用 \\
    Cookie Secure & \texttt{true} & HTTPS 生产环境必须开启 \\
    防火墙/安全组 & 22（受限）、80、443 & 关闭 8082、MySQL 和 Redis 公网暴露 \\
    \bottomrule
  \end{tabularx}
\end{table}

\section{CodeArts 流水线}

CodeArts Repo 使用 \texttt{shixun-2/master}。已归档的流水线 \#12 对提交 \texttt{41ade53e} 执行构建、代码检查、制品和部署，2026年7月14日全流程通过，总耗时 5 分 1 秒；部署目标为 ECS \texttt{1.14.150.130}，演示端口为 9092。流水线最终通过公网健康检查判定发布结果。

\begin{table}[H]
  \centering
  \caption{CodeArts 流水线阶段}
  \begin{tabularx}{\textwidth}{p{2.8cm}p{3.6cm}p{2.2cm}L}
    \toprule
    阶段 & 任务 & 结果 & 耗时 \\
    \midrule
    构建和检查 & 构建 & 成功 & 2 分 22 秒 \\
    构建和检查 & 代码检查 & 成功 & 1 分 21 秒 \\
    部署和测试 & 部署 & 成功 & 1 分 41 秒 \\
    接口测试 & Shell 健康检查 & 成功 & 56 秒 \\
    \bottomrule
  \end{tabularx}
\end{table}

\thesisfigure{华为CICD.png}{CodeArts 持续集成与持续交付概览}
\thesisfigure{华为云编译构建.png}{CodeArts 编译构建任务}
\thesisfigure{华为云制品仓库.png}{CodeArts 制品仓库}
\thesisfigure{华为云部署.png}{CodeArts 部署任务}

\chapter{项目管理与协作}

\section{代码仓库与分支协作}

项目使用 Git 进行版本管理，以 GitHub \texttt{main} 和 CodeArts \texttt{master} 作为双仓库主分支。开发变更应在独立分支形成小粒度提交，经测试和审查后合并。提交身份使用已绑定邮箱，保证贡献者统计准确；任何双仓库同步都不得使用强制推送覆盖他人历史。

\thesisfigure{GitHub仓库.png}{GitHub 项目仓库}
\thesisfigure{GitHubPR.png}{GitHub Pull Request 协作记录}
\thesisfigure{华为云代码仓库1.png}{CodeArts 代码仓库概览（一）}
\thesisfigure{华为云代码仓库2.png}{CodeArts 代码仓库概览（二）}

\section{工作项与过程管理}

项目在 CodeArts 中维护工作项和仪表盘，将需求、开发、测试、部署和缺陷闭环关联。工作项需要明确责任人、优先级、验收标准和状态；代码提交与工作项建立追踪关系，流水线结果和测试证据随版本归档。

\thesisfigure{华为云仪表盘.png}{CodeArts 项目仪表盘}
\thesisfigure{华为云工作项1.png}{CodeArts 工作项（一）}
\thesisfigure{华为云工作项2.png}{CodeArts 工作项（二）}
\thesisfigure{华为云工作项3.png}{CodeArts 工作项（三）}
\thesisfigure{华为云工作项4.png}{CodeArts 工作项（四）}

\section{质量保证与风险控制}

\begin{table}[H]
  \centering
  \caption{主要风险与控制措施}
  \begin{tabularx}{\textwidth}{p{3.2cm}p{4.8cm}L}
    \toprule
    风险 & 影响 & 控制措施 \\
    \midrule
    生产密钥泄露 & 数据库、Redis 或云资源被未授权访问 & 使用 Secret 和环境变量，不提交真实密码，定期轮换 \\
    权限绕过 & 医生获得管理员写权限 & 服务端角色校验、越权测试、权限树只作界面辅助 \\
    图片路径穿越 & 任意文件被覆盖或读取 & 路径规范化、随机文件名、目录边界校验和类型限制 \\
    数据迁移失败 & 部署中断或数据不一致 & 迁移前备份、版本化脚本、预发布验证和回滚方案 \\
    双库历史分叉 & GitHub 与 CodeArts 版本不一致 & 同一提交同步、合并前拉取、禁止强推并记录失败原因 \\
    覆盖率误读 & 以行覆盖率代替真实业务质量 & 同时报告分支、函数、API 和端到端测试结果 \\
    \bottomrule
  \end{tabularx}
\end{table}

\chapter{总结与展望}

\section{项目总结}

本项目完成了从需求分析、UML 建模、数据库设计、前后端实现、自动化测试到容器部署和持续集成的完整软件工程闭环。系统形成管理员与医生两类角色的访问边界，覆盖八类核心业务管理及首页、地图、图片和健康检查能力。代码级测试达到 241/241 通过，后端与前端行覆盖率达到 100\%，历史部署环境的 57 项 API 回归全部通过。

本报告严格依据仓库事实组织内容，区分当前代码级测试与历史部署回归，保留生产域名等尚未确认信息的待填写标记。全部项目效果图、UML 图和 GitHub/华为云证据图均纳入报告，便于课程验收和后续复核。

\section{后续改进}

后续工作包括：在隔离集成环境重新执行 57 项 API 回归并设为发布门禁；增加 Playwright 或 Cypress 浏览器端到端测试；提升前端函数覆盖率及前后端分支覆盖率；将 MySQL 与 Redis 迁入私网并启用 TLS；完成正式子域名、HTTPS、WAF、日志采集、指标告警和备份恢复演练；进一步优化前端分包和地图弱网体验。

\addcontentsline{toc}{chapter}{参考文献}
\begin{thebibliography}{99}
  \bibitem{springboot} VMware. Spring Boot Reference Documentation[EB/OL].
  \bibitem{springsecurity} VMware. Spring Security Reference[EB/OL].
  \bibitem{vue} Vue.js Team. Vue 3 Documentation[EB/OL].
  \bibitem{vite} Vite Team. Vite Documentation[EB/OL].
  \bibitem{docker} Docker Inc. Docker Compose Documentation[EB/OL].
  \bibitem{mysql} Oracle. MySQL 8.0 Reference Manual[EB/OL].
  \bibitem{redis} Redis Ltd. Redis Documentation[EB/OL].
  \bibitem{codearts} 华为云. CodeArts 软件开发生产线文档[EB/OL].
  \bibitem{projectrepo} 慧医数字医疗应用系统项目仓库、SQL、部署脚本和过程文档[Z]. 2026.
\end{thebibliography}

\appendix
\chapter{主要 API 清单}

\begin{longtable}{p{2.1cm}p{4.8cm}p{2.4cm}p{4.2cm}}
  \caption{认证与公共 API}\label{tab:public-api}\\
  \toprule
  方法 & 路径 & 权限 & 说明 \\
  \midrule
  \endfirsthead
  \toprule
  方法 & 路径 & 权限 & 说明 \\
  \midrule
  \endhead
  POST & \texttt{/api/login} & 公开 & 表单或 JSON 登录并设置 httpOnly Cookie \\
  GET & \texttt{/api/session} & 公开 & 探测并恢复当前会话 \\
  GET & \texttt{/api/permissions} & 已登录 & 按会话 \texttt{accountId} 返回 MENU 树、ACTION 权限码和角色并集，忽略客户端角色参数 \\
  POST & \texttt{/api/logout} & 已登录 & 删除 Redis 会话并清除 Cookie \\
  GET & \texttt{/api/dashboard} & \texttt{dashboard:read} & 返回首页统计、分布和资讯 \\
  GET & \texttt{/actuator/health} & 网关策略控制 & 返回应用、MySQL 和 Redis 状态 \\
  GET & \texttt{/api/regeo} & \texttt{sale-map:read} & 高德逆地理编码代理 \\
  \bottomrule
\end{longtable}

下表括号内给出后端方法级 \texttt{ACTION} 权限码。查询通常使用 \texttt{resource:read}，新增、修改和删除使用 \texttt{resource:write}；管理员拥有全量动作权限，医生只拥有只读权限。前端隐藏写按钮只是交互优化，所有受保护接口仍由 \texttt{@PreAuthorize} 独立校验。

\begin{longtable}{p{2.5cm}p{5.2cm}p{5.8cm}}
  \caption{业务 API 概览}\label{tab:business-api}\\
  \toprule
  模块 & 查询接口（ACTION） & 写操作接口（ACTION） \\
  \midrule
  \endfirsthead
  \toprule
  模块 & 查询接口（ACTION） & 写操作接口（ACTION） \\
  \midrule
  \endhead
  城市 & \texttt{GET /api/citys}、分页查询（\texttt{city:read}） & \texttt{POST /api/citys}、\texttt{DELETE /api/citys/\{id\}}（\texttt{city:write}） \\
  公司 & \texttt{GET /api/companys}、分页查询（\texttt{company:read}） & \texttt{POST}、\texttt{PUT/DELETE /api/companys/\{id\}}（\texttt{company:write}） \\
  销售地点 & \texttt{GET /api/sales}、分页查询（\texttt{sale:read}） & \texttt{POST}、\texttt{PUT/DELETE /api/sales/\{id\}}（\texttt{sale:write}） \\
  销售地图 & \texttt{GET /api/regeo}（\texttt{sale-map:read}） & -- \\
  公司政策 & \texttt{GET /api/company\_policys}（\texttt{company-policy:read}） & \texttt{POST}、\texttt{PUT/DELETE}（\texttt{company-policy:write}） \\
  医保政策 & \texttt{GET /api/medical\_policys}（\texttt{medical-policy:read}） & \texttt{POST}、\texttt{PUT/DELETE}（\texttt{medical-policy:write}） \\
  医生 & \texttt{GET /api/doctors}、\texttt{/doctors/info}（\texttt{doctor:read}） & 新增/修改/删除（\texttt{doctor:write}）；密码重置（\texttt{doctor:reset-password}） \\
  药品 & \texttt{GET /api/drugs}、分页查询（\texttt{drug:read}） & \texttt{POST}、\texttt{PUT/DELETE /api/drugs/\{id\}}（\texttt{drug:write}） \\
  材料 & \texttt{GET /api/materials}（\texttt{material:read}） & \texttt{POST}、\texttt{PUT/DELETE /api/materials/\{id\}}（\texttt{material:write}） \\
  图片 & \texttt{GET /image/\{name\}}（静态资源策略） & \texttt{POST /api/base/upload}（\texttt{file:upload}） \\
  \bottomrule
\end{longtable}

\chapter{运行与部署核对表}

\begin{table}[H]
  \centering
  \caption{部署前核对表}
  \begin{tabularx}{\textwidth}{p{1.4cm}p{4.5cm}L}
    \toprule
    序号 & 检查项 & 通过标准 \\
    \midrule
    1 & JDK、Maven、Node、Docker 版本 & JDK 17；Maven 3.8+；Node 18+；Compose v2 \\
    2 & MySQL/Redis 可达与初始化 & \texttt{medicine} 已按 \texttt{medical.sql}、V2、V3、V4、V5 顺序初始化为 23 张表 \\
    3 & 私密配置 & 密码只在 Secret 或环境变量中，权限最小化 \\
    4 & 构建测试 & \texttt{mvn verify}、前端覆盖率测试与构建成功 \\
    5 & 端口与域名 & 80/443 可达；8082 不对公网；CORS 与 Cookie 一致 \\
    6 & 健康与回归 & 容器 healthy；健康检查为 UP；业务烟测通过 \\
    7 & 回滚准备 & 上一个制品、镜像标签和数据库备份可用 \\
    \bottomrule
  \end{tabularx}
\end{table}

\chapter{图片证据完整性说明}

本报告共使用“最终交付文档/图片”目录中的 47 张 PNG 图片，其中 UML 建模图 18 张、项目运行效果图 14 张、GitHub 与华为云平台证据图 15 张。图片均以嵌入方式写入 PDF，不依赖外部链接；原始图片继续保留在交付目录，便于验收时查看细节和核对来源。

\begin{table}[H]
  \centering
  \caption{图片证据分类}
  \begin{tabularx}{\textwidth}{p{4.0cm}p{2.5cm}L}
    \toprule
    分类 & 数量 & 主要用途 \\
    \midrule
    UML 建模图 & 18 & 需求、概要设计、详细设计和数据库设计 \\
    项目运行效果图 & 14 & 登录、首页、业务管理、地图和增删改查效果 \\
    GitHub/华为云证据图 & 15 & 代码仓库、PR、工作项、测试、构建、制品和部署 \\
    合计 & 47 & 完整覆盖最终交付图片目录 \\
    \bottomrule
  \end{tabularx}
\end{table}

\end{document}
